\uuid{yOJf}
\exo7id{3745}
\titre{exo7 3745}
\auteur{quercia}
\organisation{exo7}
\datecreate{2010-03-11}
\isIndication{false}
\isCorrection{false}
\chapitre{Espace euclidien, espace normé}
\sousChapitre{Projection, symétrie}
\cpge{mpsi-mp2i.espaces-prehilbertiens.c, pcsi.espaces-prehilbertiens.c, mp-mpi.endomorphismes-euclidien.b, pc.endomorphismes-euclidiens.b, psi.endomorphismes-euclidien.b, pt.isometries-euclidien.b, mp-mpi.endomorphismes-euclidien.c, pc.endomorphismes-euclidiens.a, psi.endomorphismes-euclidien.a, pt.isometries-euclidien.a, mp-mpi.endomorphismes-euclidien.f, pc.endomorphismes-euclidiens.d, psi.endomorphismes-euclidien.f, pt.isometries-euclidien.e, bcpst2.produit-scalaire.c}
\module{Algèbre}
\niveau{L2}
\difficulte{}

\contenu{
\texte{
Soit $M \in {\cal O}(n)$.
}
\begin{enumerate}
    \item \question{Montrer que $M$ est la matrice d'une symétrie orthogonale si et seulement si
    $M$ est symétrique.}
    \item \question{Dans ce cas, déterminer la base et la direction de cette symétrie en
    fonction des matrices $I+M$ et $I-M$.}
\end{enumerate}
}
